\subsection{Password Hashing and Encryption}

User passwords \mustnot{} be stored in plain text, encoded form, or a reversible format.

Passwords used for authentication \should{} be stored using a unique, random salt and a memory-hard password-hashing function such as Argon2id. The work factor \should{} be selected based on system performance and reviewed periodically as computing power increases.

The password record \should{} include the hashing algorithm, version, salt, and work-factor parameters so that passwords can be migrated to stronger parameters in the future.

Passwords \mustnot{} be hashed using fast general-purpose functions such as MD5, SHA-1, or unsalted SHA-256.

Secrets that must later be recovered, such as service passwords, API keys, and encryption keys, \must{} be encrypted using an approved authenticated
encryption method. Encryption keys \must{} be protected separately from the encrypted secrets, preferably using a key-management system, HSM, or equivalent protected storage.

\subsection{Access Control and Accountability}

Access to stored credentials \must{} follow least privilege and need-to-know principles.

Each person \should{} use an individual account. Shared administrator accounts \should{} be avoided; where they are unavoidable, their use \must{} be logged and reviewed.

Administrative access to the password manager and infrastructure credentials \should{} require MFA, preferably using a phishing-resistant authenticator.

Access \must{} be removed promptly when personnel change roles or leave the organization.

Credential access, creation, modification, export, and recovery events \should{} be logged and reviewed. Logs \mustnot{} contain passwords, recovery
secrets, private keys, or other secret values.

\subsection{Credential Lifecycle}

Credentials \must{} have an identified owner, purpose, system, and classification.

Credentials \should{} be generated by an approved password manager or secrets-management system rather than manually chosen.

Default passwords \must{} be changed before a system is placed into service.

Credentials \must{} be rotated when they are exposed, suspected of compromise, used by departing personnel, or transferred between custodians.

Unused, obsolete, and duplicate credentials \must{} be disabled or deleted. Service-account credentials \should{} be rotated automatically where practical, and their permissions \should{} be limited to the minimum required tasks.

\subsection{Password and Secret Recovery}

Recovery procedures \must{} be documented and tested periodically.

Recovery access \should{} require approval from at least two authorized people for critical systems. Recovery codes and backup keys \should{} be one-time or tightly controlled where the system supports this.

A recovery procedure \should{} specify:

\begin{itemize}
    \item Who may authorize recovery
    \item How the identity of the requester is verified
    \item Which records or keys are required
    \item How emergency access is logged
    \item How credentials are changed after recovery
    \item How exposed or superseded credentials are revoked
    \item How the incident is documented and reviewed
\end{itemize}

\subsection{Physical Password Storage}

Physical recovery records should be treated as emergency recovery.

They \shouldnot{} contain every user's password.

They \should{} be stored in a secure location and access should be limited to authorized personnel only.

They \should{} include:
\begin{itemize}
    \item Firewall/router recovery credentials
    \item Proxmox or hypervisor emergency credentials
    \item Password manager administrator recovery information
    \item Backup encryption keys
    \item VPN recovery information
    \item Identity-server emergency administrator credentials
    \item Database administrator recovery credentials
    \item Recovery codes for critical MFA accounts
    \item A current network diagram and IP/VLAN documentation
    \item Instructions for restoring the password manager and backups
\end{itemize}

\donot{} put ordinary user passwords, private SSH keys, API keys, or unencrypted database passwords in a shared filing cabinet.

Store routine credentials in a password manager with individual access and MFA. The physical copy should contain only the minimum secrets needed to recover infrastructure.

Physical recovery records \should{} be sealed in tamper-evident envelopes or containers and stored in a fire-resistant, access-controlled location.

The organization \must{} maintain an inventory of physical recovery records, including their location, responsible custodians, creation date, and last review date.

Physical recovery records \must{} be reviewed at least annually and whenever infrastructure, personnel, or recovery procedures change.

Superseded records \must{} be destroyed securely. Sensitive paper records should be cross-cut shredded, pulped, or destroyed by an approved secure disposal provider.

The physical records \must{} not be photographed, scanned, copied, or stored in an ordinary shared office file without authorization.

\subsubsection{Split custody}

For the highest risk secrets, use split custody:
\begin{itemize}
    \item Divide a recovery secret or key into two parts so that no single person can use it alone.
    \item This is useful for root recovery, backup decryption, or the password manager emergency account.
\end{itemize}

At least two authorized custodians \should{} know that the records exist, while no single custodian should have access to the complete recovery material for the highest-risk systems.